\scanpage{178}\mspage{33}%
vérifier et d'expliciter comment les ``6 opérations'' habituelles sur les faisceaux (soit sur $X$, soit sur une réunion de strates de $X$) peuvent s'exprimer dans ce dictionnaire, dans le cas disons de strates non singulières (\unsure{autrement}, on aurait des difficultés avec les complexes dualisants, dont on aimerait qu'ils appartiennent à notre catégorie%
% A pale isolated point at approximately (949,661), below the aligned ellipsis deposits; origin uncertain, no punctuation function assigned.
\raisebox{0pt}[0pt][0pt]{\parbox{0pt}{\figure{\begin{tikzpicture}[overlay]\fill[black!45] (0,-0.06) circle[radius=0.018cm];\end{tikzpicture}}}}%
\unsure{.\,.\,.\,.\,.|\ldots}) \inserted{et de revérifier le formulaire connu de ces opérations}. Mais il apparaitra probablement que, pour bien pouvoir tout transcrire, il convient, plutôt \inserted{que} de regarder un diagramme du type
\diagram{\begin{tikzcd}[ampersand replacement=\&,column sep=1.4em,row sep=1.2em]
\cdot \arrow[r] \arrow[d] \&
% Distinct off-axis point at approximately (1229,905), below/right of the right arrowhead; it is not a target object. Its origin is uncertain.
\raisebox{0pt}[0pt][0pt]{\parbox{0pt}{\figure{\begin{tikzpicture}[overlay]\fill (0.16,-0.15) circle[radius=0.018cm];\end{tikzpicture}}}} \\
{} \& {}
\end{tikzcd}}
dans la catégorie homotopique formée avec la cat.\,. des ens. semi-simpliciaux, \inserted{de} regarder la catégorie des diagrammes d'ensembles $1/2$ simpliciaux, et à partir de là passer%
% Small isolated point at approximately (1359,1333), below the latter part of passer; origin uncertain.
\raisebox{0pt}[0pt][0pt]{\parbox{0pt}{\figure{\begin{tikzpicture}[overlay]\fill (-0.35,-0.15) circle[radius=0.016cm];\end{tikzpicture}}}} à une catégorie homotopique de %
% Isolated point at approximately (1157,1368), above/left of the f of fractions; origin uncertain.
\raisebox{0pt}[0pt][0pt]{\parbox{0pt}{\figure{\begin{tikzpicture}[overlay]\fill (0,0.37) circle[radius=0.022cm];\end{tikzpicture}}}}fractions.

J'ai plus ou moins explicité dernièrement, en cogitant sur \deleted{\unsure{les|des|ces}}des fondements de ``topologie modérée'' (i.e. où on élimine d'emblée tous les phénomènes de sauvagerie), comment une stratification équisingulière%
% Cancelled interlinear parenthesis above équisingulière; the enclosing/attachment stroke reaches the destination just before the comma.
\inserted{\deleted{(d'un espace \unsure{modéré} \unsure{compact})}}, à strates non singulières disons, d'un ``espace modéré'' \inserted{compact}, donne naissance canoniquement à un diagramme \deleted{\unsure{fon|f}} d'espaces,
% archive sheet number: 178
